---
format: yohaku/0.2
target: text
macro: '[[macros/raw]]'
imported: {"bytes":6545,"path":"/root/projects/yohaku/release-prep/yohaku/agent/yohaku_codex_linux.py","sha256":"8e489ce1a4298f05d1bdafca7d35267b114f6c2a0319e41f3acbc65c8016a3eb"}
---

## 正文
```text
"""Native Linux Codex app-server runtime for Yohaku.

Host inference is separate from world tools. No Windows setup and no automatic
credential import. Native command/exec is verified in a deny-by-default policy.
"""
import json
import os
from pathlib import Path
import shutil
import stat
import subprocess
import tempfile

# Import the existing bridge module only when needed; avoid import cycles.
import yohaku_apps as apps

PROFILE='yohaku-control'
PROBE='''import json,pathlib,sys
cwd,outer=map(pathlib.Path,sys.argv[1:])
r={}
try:r['control_read']=(cwd/'read-control').read_text()=='yohaku-readable-control'
except OSError:r['control_read']=False
for key,path in [('workspace_write_denied',cwd/'write-probe'),('outside_write_denied',outer/'outside-write')]:
 try:path.write_text('probe');r[key]=False
 except OSError:r[key]=True
try:(outer/'outside-read').read_bytes();r['outside_read_denied']=False
except OSError:r['outside_read_denied']=True
print(json.dumps(r))
'''


def is_linux_codex(address):
    try:
        path=apps.address_path(address)
        if path.name!='codex' or not path.is_file() or not os.access(path,os.X_OK):return False
        with path.open('rb') as stream:hdr=stream.read(20)
        return len(hdr)==20 and hdr[:6]==b'\x7fELF\x02\x01' and hdr[18:20]==b'\x3e\x00'
    except (OSError,ValueError,apps.tools.ToolError):return False


def _auth_home():
    return Path.home()/'.local/state/yohaku/codex-linux-login'


def _runtime_config(executable,cwd):
    return '''approval_policy = "never"
default_permissions = "yohaku-control"
cli_auth_credentials_store = "file"
mcp_servers = {}
hooks = {}
plugins = {}
[features]
shell_tool = false
multi_agent = false
apps = false
plugins = false
use_legacy_landlock = false
[permissions.yohaku-control.filesystem]
":root" = "deny"
":minimal" = "read"
'''+json.dumps(str(executable)) + ' = "read"\n' + json.dumps(str(cwd))+''' = "read"
[permissions.yohaku-control.network]
enabled = false
'''


def _safe_auth_copy(target):
    # Login is explicit and credentials remain solely on the user's device.
    source=_auth_home()/'auth.json'
    try:
        metadata=source.lstat()
    except FileNotFoundError:
        return
    if not stat.S_ISREG(metadata.st_mode) or metadata.st_size>1024*1024:
        raise ValueError('Invalid dedicated Linux Codex authentication file')
    # Open without following a replaced symlink.
    fd=os.open(source,os.O_RDONLY|os.O_NOFOLLOW|os.O_CLOEXEC)
    try:
        current=os.fstat(fd)
        if not stat.S_ISREG(current.st_mode) or current.st_size>1024*1024:
            raise ValueError('Invalid Codex credential file')
        with os.fdopen(fd,'rb',closefd=False) as src:
            with target.open('xb') as dest:
                dest.write(src.read(1024*1024+1))
        target.chmod(0o600)
    finally:os.close(fd)


def launch(address):
    if not is_linux_codex(address):
        raise ValueError('Choose a native x86_64 Linux Codex ELF executable')
    executable=apps.address_path(address)
    python=shutil.which('python3',path='/usr/bin:/bin')
    if python is None:raise ValueError('System python3 is needed for the Linux boundary probe')
    cache=Path.home()/'.cache/yohaku'
    cache.mkdir(mode=0o700,parents=True,exist_ok=True)
    runtime=tempfile.TemporaryDirectory(prefix='codex-bridge-',dir=cache)
    try:
        outer=Path(runtime.name);home=outer/'state';cwd=outer/'control'
        home.mkdir(mode=0o700);cwd.mkdir(mode=0o700);(home/'tmp').mkdir(mode=0o700)
        (cwd/'read-control').write_text('yohaku-readable-control')
        (outer/'outside-read').write_text('outside-canary')
        (home/'config.toml').write_text(_runtime_config(executable,cwd))
        (home/'config.toml').chmod(0o600)
        _safe_auth_copy(home/'auth.json')
        env={'PATH':'/usr/bin:/bin:/usr/sbin:/sbin','HOME':str(home),'CODEX_HOME':str(home),
             'TMPDIR':str(home/'tmp'),'LANG':'C.UTF-8','PYTHONDONTWRITEBYTECODE':'1'}
        child=subprocess.Popen([str(executable),'app-server','--listen','stdio://'],env=env,cwd=cwd,
                               stdin=subprocess.PIPE,stdout=subprocess.PIPE,stderr=subprocess.DEVNULL,start_new_session=True)
        child.yohaku_linux=True
        child.yohaku_runtime=runtime
        child.yohaku_info={'cwd':str(cwd),'outside':str(outer),'readProbe':str(outer/'outside-read'),
                           'python':python,'executable':str(executable)}
        return child
    except BaseException:
        runtime.cleanup()
        raise


def verify_boundaries(connection):
    from yohaku_codex import CodexError
    info=connection.info
    result=connection.request('command/exec',{'command':[info['python'],'-I','-B','-c',PROBE,info['cwd'],info['outside']],
        'cwd':info['cwd'],'permissionProfile':PROFILE,'timeoutMs':10000},timeout=25)
    if result.get('exitCode')!=0:raise CodexError('Linux Codex sandbox probe failed to execute; no model turn started')
    raw=result.get('stdout')
    if not isinstance(raw,str) or len(raw)>4096:raise CodexError('Linux Codex sandbox probe returned invalid output')
    try:
        checks=json.loads(raw)
    except ValueError:raise CodexError('Linux Codex sandbox probe returned invalid JSON')
    expected={'control_read','workspace_write_denied','outside_write_denied','outside_read_denied'}
    if not isinstance(checks,dict) or set(checks)!=expected or any(type(checks[k]) is not bool for k in expected):
        raise CodexError('Linux Codex sandbox probe returned invalid checks')
    # A tmpfs write inside bwrap does not imply a write to the host filesystem.
    host_blocked=not (Path(info['outside'])/'outside-write').exists()
    if not all(checks[k] for k in ('control_read','workspace_write_denied','outside_read_denied')) or not host_blocked:
        raise CodexError('Linux Codex sandbox boundary check failed; no model turn started')


def cleanup(child):
    runtime=getattr(child,'yohaku_runtime',None)
    if runtime is not None:runtime.cleanup()


def login(address):
    """Only performed by an explicit user-invoked login command; never during probe."""
    if not is_linux_codex(address):raise ValueError('Choose a native Linux Codex executable')
    directory=_auth_home();directory.mkdir(mode=0o700,parents=True,exist_ok=True)
    env=os.environ.copy();env['CODEX_HOME']=str(directory)
    print('Sign in to dedicated Yohaku Linux Codex storage; no Windows credentials are copied.',flush=True)
    return subprocess.call([str(apps.address_path(address)),'login','--device-auth'],env=env)
```

## 输出与关联
```yaml
output: "agent/yohaku_codex_linux.py"
```
